\documentclass[a4paper,12pt]{article}

\usepackage[utf8]{inputenc} 
\usepackage{amsmath}
\usepackage[a4paper, margin=0.8in, top=0.1in]{geometry} % Adjust the margin here
\usepackage{multirow}
\usepackage{nopageno}
\usepackage{graphicx} % Add this line to include graphics
\usepackage{wrapfig}
\usepackage{tikz}

\setlength{\parindent}{0pt} % Set global no indent
    
\title{{\large Department of Physics, FNSPE, CTU in Prague} \\
                  02YMECH - Mechanics \\ 
          {\Large \textbf{Midterm 2}}\vspace{-1cm}}
\date{19.12.2025 / 14:00 - 15:30}
\author{}

\begin{document}

\maketitle

\textbf{Q1:} Consider a system of three point particles, each having an identical mass $m$, with their positions defined by time-dependent vectors,
\begin{equation*}
    \vec r_{1}~=~(at^2,~vt,~0),~\vec r_{2}~=~(vt,~at^2,~b),~\vec r_{3}~=~(0,~b,~at^2), 
\end{equation*}
where $a$, $b$, and $v$ are positive constants.
\begin{itemize}
    \item[\textbf{a.}] Find the total force required to produce the motion. --- \textit{0.2 points}
    \item[\textbf{b.}] Find the total angular momentum and torque concerning the origin of the coordinate system. --- \textit{0.3 points}
    \item[\textbf{c.}] Based on the results from parts (a) and (b), determine if total linear momentum and total angular momentum are conserved. --- \textit{0.2 points}
    \item[\textbf{d.}] Find the total kinetic energy in the LAB frame and in the center of mass frame. --- \textit{0.3 points}
\end{itemize}

\vspace*{0.2cm}
\textbf{Q2:} The equation of motion of a damped oscillator is given as $x(t)=e^{-\beta t}x_0\cos{\left(\omega_1 t\right)}$. How long does it take for the damped oscillator to halve the amplitude of the oscillation? 
(Givens: $\ln2\approx 0.7$, $m=~1\text{kg}$, $b=1.4~\text{kg/s}$) --- \textit{1 point}

\vspace*{0.2cm}
\textbf{Q3:} For a given central force $f(r)$ the orbit of a particle of mass $\mu$ can be calculated by integrating the following differential equation,
\begin{equation*}
    \dfrac{d^2}{d\theta^2}\left(\frac{1}{r}\right)+\frac{1}{r}=-\frac{\mu}{l^2}r^2f(r)
\end{equation*}
where $l$ is the angular momentum of the particle. The particle follows an orbit given by $1+\epsilon\cos\theta=\alpha/r$, where $\alpha\equiv \frac{l^2}{\mu k}$ and $\epsilon \equiv \sqrt{1+\frac{2El^2}{\mu k^2}}$ are constants.

\begin{itemize}
    \item[\textbf{a.}] Find the force law $f(r)$ for the central force in terms of $k$ and $r$. --- \textit{0.6 points}
    \item[\textbf{b.}] Find the potential $U(r)$ associated with $f(r)$. --- \textit{0.4 points}
\end{itemize}

\vspace*{0.2cm}
\textbf{Q4:} The earth rotates about its axis with angular velocity $\vec{\omega}$. The unit vectors $\hat{e}_1$, $\hat{e}_2$, and $\hat{e}_3$ are local unit vectors on the Earth's surface, where $\hat{e}_3$ is perpendicular to the surface, while $\hat{e}_1$ and $\hat{e}_2$ lie on the surface. The component of $\vec{\omega}$ along the $\hat{e}_3$ direction is $\omega_z$. A very fast particle of mass $m$ is projected on a horizontal surface with an initial velocity $\vec{v_0}=v_0\hat{e}_2$. Assume the Earth is a perfect sphere with radius $R$ and neglect air resistance.

\begin{itemize}
    \item[\textbf{a.}] Determine the Coriolis force by considering the general velocity of the particle $\vec{v}=v_1\hat{e}_1+v_2\hat{e}_2$. Describe qualitatively the deflection due to the Coriolis force. --- \textit{0.6 points}
    \item[\textbf{b.}] Write the coupled equations of motion for $\hat{e}_1$ and $\hat{e}_2$ directions, and solve them for $v_1$ and $v_2$. Does this result satisfy your initial expectations for the trajectory? --- \textit{0.4 points}
\end{itemize}


(Given: Coriolis force $\vec{F}_c = -2m \vec{\Omega} \times \vec{v}$)

\end{document}

